% answers_short.tex —— 2008 年·普通高等学校招生考试·辽宁卷·文科 · 简短答案速查
% 由 tools/make_short_answers.py 自 answers.tex 的 \gksolhead/\gkanswer 提取，
% 并与 2008/2008_数学真题_辽宁文_答案要点.md 逐题比对后生成，勿手改（改源 MD 后重跑）。

\begin{gknotice}
\textbf{使用说明}：本节只给最终答案便于快速核对，完整推导见解析版。答案为编写者独立推导，\textbf{非官方标准答案}。
\end{gknotice}

\gkpart{选择题}
\begin{center}\begin{tabular}{ccccccccccccc}
\toprule
题号 & 1 & 2 & 3 & 4 & 5 & 6 & 7 & 8 & 9 & 10 & 11 & 12 \\
\midrule
答案 & D & C & B & C & A & A & C & A & B & B & D & D \\
\bottomrule
\end{tabular}\end{center}

\gkpart{填空题}
\begin{itemize}[leftmargin=2.6em]
  \item[\textbf{13.}] $y=\dfrac12(\ln x-1)$（$x>0$）。
  \item[\textbf{14.}] $\dfrac32$。
  \item[\textbf{15.}] $35$。
  \item[\textbf{16.}] $\sqrt3$。
\end{itemize}

\gkpart{解答题}
\begin{itemize}[leftmargin=2.6em]
  \item[\textbf{17.}] （1）$a=b=2$；（2）$\dfrac{2\sqrt3}3$。
  \item[\textbf{18.}] （1）$0.2$，$0.5$，$0.3$；（2）① $0.7599$；② $0.0621$。
  \item[\textbf{19.}] （1）证明见解析；（2）定值 $\sqrt2$；（3）$\dfrac{\sqrt2}2$。
  \item[\textbf{20.}] （1）是等比数列；（2）前 $n$ 项和为 $4n$。
  \item[\textbf{21.}] （1）$x^2+\dfrac{y^2}4=1$；（2）$k=\pm\dfrac12$，$\left|AB\right|=\dfrac{4\sqrt{65}}{17}$。
  \item[\textbf{22.}] （1）$b=0$；增区间 $(-\infty,-1)$、$(1,+\infty)$，减区间 $(-1,1)$。（2）$\left[-\dfrac{2\sqrt3}3,\dfrac{2\sqrt3}3\right]$。
\end{itemize}
